\section{Radial regeneration and normalized correlations}\label{sec:radial}

The first normalization problem is local. A loop surrounding a marked
point can be arbitrarily small, so convergence of macroscopic loops does
not directly give convergence of an insertion which weights every such
loop. We prove that the dependence on these small loops cancels in ratios
of partition functions. The cancellation must remain bounded under a
complex perturbation of the weights, because Section~\ref{sec:electric}
will use analytic continuation.

We work in unit tile coordinates. For a rounded loop $\ell$ and a tile
vertex $u$, put
\[
 R_u(\ell)=\max_{z\in\ell}|z-u|.
\]
Enclosure refers to the disjoint switch drawing from
Section~\ref{sec:inputs}. A circuit state $\gamma$ consists of an actual
primal-open circuit surrounding $u$, with all its vertices wired, and
the ordinary FK law in its interior. Its size is $R$ if its inscribed
and outer radii about $u$ lie between fixed positive multiples of $R$.
No further regularity of the circuit is required. At a circuit edge,
the two switch arcs belong to its respective sides. Consequently the
interior loop configuration separates exactly from the exterior, also
when a loop uses an inward arc beside the circuit.

Let $N_\gamma$ count its interior loops enclosing $u$, and define
\begin{equation}\label{eq:radial-partition}
 F(\gamma;v)=\E_\gamma v^{N_\gamma},
 \qquad
 \frac{\dd\Pp_{\gamma,v}}{\dd\Pp_\gamma}
       =\frac{v^{N_\gamma}}{F(\gamma;v)},\qquad v>0.
\end{equation}
A plane state of size $R$ uses the plane FK law and counts precisely
the loops enclosing $u$ with $R_u(\ell)\le R$. We use the notation
$F(\gamma;v)$ for this state as well. An exchange of the two colors
gives the corresponding statements for dual circuits. In comparisons
the color of $u$ is fixed, or transported by a lattice symmetry together
with the two color names.

\begin{proposition}[Radial estimates]\label{prop:radial}
Fix $1\le q<4$, fixed radius-comparison constants, and a compact interval
$V\subset(0,\infty)$. The following bounds are uniform over the above
circuit and plane states and over $v\in V$.

For states $\gamma,\gamma'$ of sizes $R,R'\ge r_0$,
\begin{equation}\label{eq:radial-comparison}
 \left|\log\frac{F(\gamma;v)}{F(\gamma';v)}\right|
       \le C\bigl(1+|\log(R/R')|\bigr).
\end{equation}
There is $\delta>0$ such that every $F(\gamma;ve^w)$ is nonzero for
$|w|<\delta$. Its logarithm, continued from its real value at $w=0$,
has absolute value at most $C\log(2+R)$. Moreover
\begin{equation}\label{eq:radial-complex-ratio}
 \left|\log\frac{F(\gamma;ve^w)}{F(\gamma;v)}
       -\log\frac{F(\gamma';ve^w)}{F(\gamma';v)}\right|
 \le C|w|\bigl(1+|\log(R/R')|\bigr).
\end{equation}

For two states of comparable size $R$, their positively tilted laws
can be coupled so that their configurations agree inside a common
surrounding primal-open circuit containing $B(u,s)$, except with
probability $C(s/R)^c$, for $r_0\le s\le R/C$. Conditional on the common
circuit, the common continuation has its circuit law tilted by $v$.

Finally, let $F_{\ge h}(\gamma;v)$ count only the eligible loops with
$R_u(\ell)\ge h$. For comparable starts of size $R$,
\begin{equation}\label{eq:radial-truncation}
 \frac{F_{\ge\epsilon R}(\gamma;v)/
             F_{\ge\epsilon R}(\gamma';v)}
      {F(\gamma;v)/F(\gamma';v)}=1+o(1)
\end{equation}
uniformly as $\epsilon\downarrow0$ and $\epsilon R\to\infty$.
The same real comparison and regeneration bounds hold with any such
lower cutoff on the counted radii.
\end{proposition}

\subsection{Passage bounds at a circuit and under a tilt}

The FK estimates used here are the bounded-density collar lemma,
exponential passage moments, and local tilt bounds in
\cite[Section~2, ``Comparison through a collar'' and ``Passages and
positive local tilts'']{FKCompanion}; their precise scope is recorded
in Input~\ref{inp:fk}. We first explain the extension to an irregular
wired circuit, where a bulk-buffer assertion alone would not suffice.

A newly counted enclosing loop whose maximum radius crosses a specified
cutoff has a passage in one of a fixed number of buffers at that scale:
it either goes around $u$ in the band or leaves the band. The same
observation charges every enclosing loop whose status changes when a
subannulus is resampled. Loops which do not enclose $u$ need not be
charged. Thus, for passage counts $X_j$ on $l$ geometrically spaced
bands, finite coloring of their buffers and conditional iteration give
\begin{equation}\label{eq:radial-band-moments}
 \E\exp\left(t\sum_{j=1}^l X_j\right)
       \le \exp(C_t(l+1)),\qquad t<\infty.
\end{equation}
The bound is conditional on data beyond the buffered cuts.

At the outer wall of a primal-wired circuit, count large enclosing
loops using their dual sides. Those sides have free boundary at the
wall. In a cropped annular or rectangular passage region, the ordered
medial separators force distinct actual dual crossing components:
two neighboring sides may use one component, so a factor two and
bounded end effects are allowed. The crop ends lie strictly inside
the region traversed by the separators. An inward switch arc beside
an open circuit edge still has an actual interior dual gap on its
other side; it cannot use a fictitious path through the exterior.

Explore complete actual dual components from one end of the crop,
including unsuccessful components. After a component is removed, its
incident unexplored boundary edges are closed in the explored color.
The remaining law is dominated by filling the cropped box and wiring
its cut boundary. The original wall is free in this color, and
conditioning farther away can identify cut-boundary vertices only.
Filling the missing vertices and, if necessary, wiring the entire
boundary can only increase actual crossing probabilities. A blocker
in the filled box has uniformly positive probability by the crossing
bounds of Input~\ref{inp:fk}, including the wall version of
\cite[Theorem~1.2]{DMT}. Successive component discoveries therefore
give a geometric tail. To obtain an arbitrary fixed exponential rate,
subdivide the passage distance into $J$ separated slices, cover each
by buffered small boxes, and apply the same exploration to the boxes
intersected with the circuit interior. Every original passage crosses
a small buffer in each slice. If $X$ is the original count and $X_i$
the $O(J^2)$ small-buffer counts, then $JX\le\sum_iX_i$.
Their buffers have a coloring with $k$ colors independent of $J$.
Conditional iteration in each color class, followed by H\"older
across classes, therefore bounds
$\E e^{cJX/k}$ using the common first exponential rate $c$ of
the $X_i$. The prefactor may depend on $J$, which is harmless:
choose $J$ large and fixed to obtain any prescribed rate.
The finitely many scales where these boxes do not fit are covered by
a deterministic finite-size bound. This proves
\eqref{eq:radial-band-moments} also in the outer circuit collar.
The argument uses the specified wired wall, not an arbitrary partition
on a jagged boundary.

Now first count nests only up to a radius a fixed factor inside the
inscribed radius of the circuit. The local tilt bounds compare
successive cutoff masses and give every fixed passage moment at the
last weighted level. They are uniform for $v\in V$: in their proof,
the H\"older exponents can be chosen using
$\sup_{v\in V}|\log v|$, on either side of $1$. Moving the cutoff
back a bounded number of levels leaves the passage buffers outside
the tilt. H\"older's inequality then restores those levels. The
preceding wall estimate restores the remaining boundedly many outer
levels at bounded cost. Negative powers have the same comparison,
by Jensen's inequality followed, when needed, by this cutoff argument.
Bounded densities across a full bulk collar now compare the interior
cutoff mass to that for any other comparable state. This proves
\eqref{eq:radial-comparison}, first for comparable sizes and then
by multiplying comparisons across scales.

These statements remain uniform if counting begins at a lower
maximum-radius threshold $h$. Start the cutoff induction with no
weight below $h$. In particular, under a tilt supported below the
inner end of $l$ bands, their joint passage moments still obey
\eqref{eq:radial-band-moments}; the innermost boundedly many buffers
use the local tilt bounds, and all others are separated from its
support. Adding the weights in those $l$ bands changes the density,
including its normalization, by a factor whose fixed moments cost
$\exp(C(l+1))$. Hence \eqref{eq:radial-band-moments} holds under the
full positive tilt up to the outer end of those bands as well.
These are the real estimates needed for
regeneration.

\subsection{Stopping circuits and exponential waiting times}

Fix $K$ sufficiently large and use levels $r_i=r_0K^i$, $i\ge0$,
with $r_0$ large enough for all crossing constructions. Each trial
band has radii comparable to $r_i$ and $3r_i$ and has full collars
on both sides. The value of $K$ leaves these enlarged trial regions
disjoint and leaves an unused inward collar after a starting circuit.
Assign a state the level just below its inscribed radius; states at
one level have sizes within a fixed factor. A plane cutoff is assigned
in the same way, and its first trial lies well inside that cutoff.

Inspect trials inward until finding a surrounding primal-open
circuit, choosing the outermost circuit in that band. The usual
outer exploration is a stopping exploration: no edge strictly inside
the selected circuit need be queried. Intersecting candidate circuits
can be merged by taking their outer boundary. Failed tests concern
previous, disjoint bands. At level zero use a terminal symbol with
continuation mass one. Conditional on a stopped circuit, the ordinary
continuation is its wired FK law. If $D$ counts the weighted enclosing
loops outside that circuit, then
\begin{equation}\label{eq:radial-transition}
 F(\gamma;v)=\E_\gamma\!\left[v^D F(\gamma_1;v)\right],
\end{equation}
where $\gamma_1$ denotes the stopped state, including the terminal
case. For a plane start, $D$ also respects its original maximum-radius
cutoff. The assignment of circuit-side arcs makes this factorization
exact. It follows that the stopped process remains a regeneration
process under the tilt.

Let $J$ be the number of levels crossed in one transition. Assigning
the loops counted by $D$ to their maximum-radius bands gives, under
any of the positive tilts under discussion,
\begin{equation}\label{eq:radial-jump-cost}
 \E_{\gamma,v}[e^{tD};J=j]\le e^{C_t(j+1)},\qquad t<\infty.
\end{equation}
We next prove
\begin{equation}\label{eq:radial-jump-tail}
 \Pp_{\gamma,v}(J>l)\le Ce^{-cl}.
\end{equation}
This needs more than a crossing bound after arbitrary neighboring
edges have been exposed, because such an exposure need not leave a
favorable conditional law under the tilt.

For each of the first $l$ trials, designate its central band as a
region to be resampled by the ordinary FK conditional law. All counts
used below have buffers in that trial's enlarged collar. A trial is
called typical if the old passage counts which can charge changed
enclosing loops are at most $H$, and, for every boundary partition
which changes outside its collar, the ordinary conditional redraw
has probability at least $b$ of producing a surrounding circuit and
at most $H$ new charging passages. Here $b>0$ and $H<\infty$ will be
fixed independently of $l$ and of the lattice scale. The second
condition depends only on collar edges outside the resampled band,
so remains true after other trial bands are changed.

We justify that typicality fails with arbitrarily small ordinary
conditional probability. Expose the collar up to its far cuts, and
let $Y$ count its actual components joining the band-facing cuts to
the far cuts, on both radial sides, plus one. Only these components
can acquire new identifications when the far partition is changed.
For fixed collar edges, any two redraw kernels thus have density
ratio at most $q^{C Y}$. Complete-component exploration in cropped
subcollars, as above, gives a uniform exponential tail for $Y$.
First fix a large bound $Y_0$. Given $Y\le Y_0$, the worst conditional
probability, over the induced far partitions, of more than $H$ new
passages is at most $q^{CY_0}$ times that probability under the
actual partition, or its conditional mixture. The latter probability
has arbitrarily small expectation as $H\to\infty$, by the ordinary
conditional passage estimate. Markov's inequality then makes the
collar-edge event on which this worst probability exceeds a fixed
tolerance arbitrarily rare. Success without the passage restriction
has conditional probability at least $2b>0$ by RSW. Choose the
tolerance below $b$, and then $H$ large. The old passage test is
handled by the same moment bound.

All these estimates hold conditional on data outside the enlarged
trial collar, uniformly in its induced partition. For any prescribed
set of trials their failures therefore have probability at most a
chosen small constant to the size of that set, by conditional
iteration. A binomial sum proves, at any desired exponential rate,
that fewer than $l/3$ typical trials is unlikely under ordinary FK.
To transfer this statement to the fully tilted law, first retain
only the tilt below the inner end of these $l$ trials, omitting
boundedly many end trials for the buffers. The same ordinary
conditional argument still applies: the retained weight counts
loops wholly inside that lower cutoff, so is measurable outside
each enlarged trial collar. Conditional on the exterior of a collar,
it is constant, and the conditional law in the collar is ordinary
FK. The remaining change of density
has second moment at most $e^{C(l+1)}$, by
\eqref{eq:radial-band-moments} and the real mass comparisons.
Cauchy--Schwarz, choosing the preceding ordinary rate larger than
$C$, gives
\begin{equation}\label{eq:radial-typical}
 \Pp_{\gamma,v}(|T|<l/3)\le Ce^{-cl},
\end{equation}
where $T$ is the set of typical trials. This proof applies equally
when a lower radius threshold is present.

For completeness, the insertion argument converting
\eqref{eq:radial-typical} into \eqref{eq:radial-jump-tail} is as
follows. Write $A_l$ for failure at all $l$ trials, and for a fixed
subset $S$ of the trials put $A_S=A_l\cap\{S\subset T\}$. Redraw
the bands in $S$ successively with their ordinary FK heat baths,
retaining only outcomes with success and at most $H$ new passages.
Each retention has probability at least $b$, even after the preceding
redraws. At one accepted redraw, at most $CH$ old or new enclosing
loops change their weight. Consequently the ratio of input to output
tilt densities is at most
$M=\exp(CH\sup_{v\in V}|\log v|)$ per redraw.
Ordinary heat baths preserve the ordinary law, including their
successive composition. If $E_S$ is the event with exact success
set $S$ among these $l$ trials, integration of the accepted transitions
therefore gives
\[
 \Pp_{\gamma,v}(E_S)
       \ge (b/M)^{|S|}\Pp_{\gamma,v}(A_S).
\]
The events $E_S$ are disjoint. With $\delta=b/M$, summing over $S$
yields the useful integrated inequality
\begin{equation}\label{eq:radial-insertion}
 1\ge
 \E_{\gamma,v}\left[\1_{A_l}(1+\delta)^{|T|}\right].
\end{equation}
The set $T$ stays inside the expectation; no deterministic choice of
typical trials is required. Equations~\eqref{eq:radial-typical} and
\eqref{eq:radial-insertion} prove \eqref{eq:radial-jump-tail}.

Combining the tail with \eqref{eq:radial-jump-cost} gives a joint
exponential moment: there exist $a,t>0$ such that
\begin{equation}\label{eq:radial-joint-moment}
 \sup_{\gamma,v}\E_{\gamma,v}e^{aJ+tD}<\infty.
\end{equation}
Indeed fix an exponent $T_0>0$ in
\eqref{eq:radial-jump-cost} and apply H\"older on $\{J=j\}$:
\[
 \E[e^{tD};J=j]
 \le \E[e^{T_0D};J=j]^{t/T_0}
             \Pp(J=j)^{1-t/T_0}.
\]
For $t/T_0$ small, the resulting exponent in $j$ is negative; then
choose $a>0$ smaller than its absolute value.

\subsection{Coupling states at different levels}

At a common level, the next-band transitions of any two states have
a common part of uniformly positive mass. To see this, use bounded
densities through the unused collar to minorize both ordinary
band marginals by a fixed fraction of one reference marginal, and
restrict to its successful circuits. Given the next screen, its
continuation mass compares to any reference mass at that level by
\eqref{eq:radial-comparison}. The number of newly charged loops can
be restricted to $D\le H_1$ at arbitrarily small discarded ordinary
probability, by the fixed-band passage estimate. On this restriction
$v^D$ is bounded above and below uniformly. Subtract the two discarded
parts from the reference marginal and then take its common minimum.
The normalizations also compare to the same continuation mass, so
the common part retains a fixed mass $p_0>0$ under the tilted laws.
In particular a jump of one level has probability bounded below.

For starts at the same or neighboring levels, advance whichever
state has the higher remaining level; when the levels agree, attempt
coupling on this common part. Use a fixed mass strictly smaller than
the available overlap. After a failed attempt the residual transition
moments are bounded by the original moments divided by $1-p_0$.
This avoids conditioning estimates depending on an overlap which
might itself be close to one.

Here is the quantitative reason this procedure also works for
asynchronous levels. If their current difference is $d\ge1$ and
the higher state jumps by $J$, the new difference is $|d-J|$.
For sufficiently small $a>0$, \eqref{eq:radial-joint-moment} gives
an exponential drift outside a fixed finite set:
\[
 \E e^{a|d-J|}
 \le e^{ad}\E[e^{-aJ};J\le d]
          +\E[e^{a(J-d)};J>d]
 \le \theta e^{ad},\qquad d\ge d_0,
\]
with $\theta<1$. The last tail term is exponentially small relative
to $e^{ad}$, while $J\ge1$ bounds the first coefficient below one.
From $\{0,\ldots,d_0\}$, finitely many one-level jumps bring the
levels together with uniformly positive probability, followed by
the common-part coupling. The exponential drift bounds return
times, and the failed-attempt residual moments remain bounded.
Geometrically summing these returns proves an exponential moment
for the number of transitions until meeting. More explicitly, the
kernel of one unsuccessful excursion from this finite set has mass
strictly below one at discount one and a bounded exponential moment;
the same is true at a sufficiently small positive time exponent.
Summing its powers proves the assertion.

Iterating \eqref{eq:radial-joint-moment} and reducing its exponents
then bounds total distance and charged-loop cost up to meeting. If
$L$ denotes the largest number of levels descended by either chain,
and $D_1^*,D_2^*$ their respective costs before the common state,
\begin{equation}\label{eq:radial-meeting-moment}
 \E\exp\{c(L+D_1^*+D_2^*)\}\le C.
\end{equation}
One may require each chain to take one transition before permitting
meeting; the bound is unchanged. The terminal symbol provides the
same construction at the finitely many bottom levels.

If the initial size is $R$, reaching a radius comparable to $s$
without meeting requires order $\log(R/s)$ descended levels.
Equation~\eqref{eq:radial-meeting-moment}, with a smaller exponent
to cover a final overshoot, bounds this probability by $C(s/R)^c$.
On meeting, both chains have exactly the same surrounding circuit
and the same conditional tilted continuation. Their configurations
can therefore be made identical inside it. This proves the inward
coupling assertion in Proposition~\ref{prop:radial}.
Terminal exhaustion, though a valid meeting for partition-function
identities, is counted as failure in this spatial coupling; the
coupling attempts here are stopped before the radii drop below $s$.

\subsection{Complex stability and removal of the inner cutoff}

Fix a real $v\in V$. For a state $\gamma$ at level $i$ define
\[
 G_\gamma(w)
     =\frac{F(\gamma;ve^w)}{F(\gamma;v)}
     =\E_{\gamma,v}e^{wN_\gamma}.
\]
At finite scale these are holomorphic functions. We prove by induction
on the level that they do not vanish and that, for states at the same
or neighboring levels,
\begin{equation}\label{eq:radial-induction}
 |\log G_\gamma(w)-\log G_{\gamma'}(w)|\le B|w|.
\end{equation}
All logarithms start at zero at $w=0$. The terminal levels have
uniformly bounded size and give the initial assertion directly.

To compare two new starts, use the preceding coupling, requiring one
transition in each chain so that their meeting state lies below the
highest level being proved. Regeneration gives
\[
 G_{\gamma_j}(w)=
        \E[e^{wD_j^*}G_{\gamma_*}(w)],\qquad j=1,2,
\]
with the same random meeting state $\gamma_*$. Divide both sides by
$G_\sigma(w)$ for any reference state $\sigma$ at the preceding
level. Chaining the inductive bounds through the lower levels gives
\[
 Q(w)=\frac{G_{\gamma_*}(w)}{G_\sigma(w)},\qquad
 |\log Q(w)|\le B|w|(L+2).
\]
By \eqref{eq:radial-meeting-moment}, for each fixed $B$ and all
sufficiently small $|w|$, both divided expectations lie in
$\{z:|z-1|<1/4\}$. Their difference has the sharper bound
\begin{align*}
 \left|\E[(e^{wD_1^*}-e^{wD_2^*})Q(w)]\right|
 &\le |w|\E\left[(D_1^*+D_2^*)
   e^{|w|(D_1^*+D_2^*)+B|w|(L+2)}\right]\\
 &\le C_0|w|.
\end{align*}
Here $C_0$ is independent of $B$, provided $B|w|$ is smaller than
a fixed fraction of the exponent in
\eqref{eq:radial-meeting-moment}. The logarithm is uniformly
Lipschitz on the displayed disk. Thus choose $B$ larger than the
resulting fixed multiple of $C_0$, and only then choose $\delta>0$
so that the disk and exponent conditions hold for $|w|<\delta$.
This closes both the nonvanishing induction and
\eqref{eq:radial-induction}. Telescoping levels proves
\eqref{eq:radial-complex-ratio}. The real comparison to a terminal
state bounds $|\log F(\gamma;v)|$ by $C\log(2+R)$, and another
telescoping sum proves the asserted logarithmic bound in the complex
neighborhood.

It remains to compare full and thresholded masses. Put $h=\epsilon R$
and let $l$ be the number of levels between $R$ and $h$. In each
regeneration path stop at its first screen whose level is at most
$J$ levels above $h$, where $1\ll J\le l/2$. Call this stop typical
if its level is between one and $J+O(1)$ levels above $h$.
The probability of an atypical stop is $Ce^{-cJ}$ under both the
full and thresholded tilts: sum the jump-tail bound over its possible
origins, each level being visited at most once. On a typical path
every loop already charged encloses the stopping screen or lies
outside it, and hence has maximum radius above $h$. The two
weighting rules have therefore agreed up to this stop.

Write
\[
 Q_h(\sigma)=F_{\ge h}(\sigma;v)/F(\sigma;v)
\]
for the continuation quotient. Normalize by its value $q_h$ at any
reference screen at the threshold level. Real comparisons, including
their thresholded version, give on typical stops
\begin{equation}\label{eq:radial-cutoff-quotient}
 e^{-C(J+1)}\le Q_h(\sigma)/q_h\le e^{C(J+1)}.
\end{equation}
Couple the two full-tilt paths from the original comparable states.
Their stopped screens agree except with probability $Ce^{-cl}$,
after changing $c$ to account for $J\le l/2$. On typical stopped
paths regeneration gives the exact identity
\[
 I_\gamma:=\E_{\gamma,v}
       [\1_{\mathrm{typ}}Q_h(\sigma)]
   =\frac{F_{\ge h}(\gamma;v)}{F(\gamma;v)}
            \Pp_{\gamma,v,\ge h}(\mathrm{typ}),
\]
where the probability on the right uses the thresholded tilt.
The difference of the two $I_\gamma$, divided by $q_h$, is
at most $Ce^{C(J+1)-cl}$. The lower bound in
\eqref{eq:radial-cutoff-quotient} converts this into a relative
error $Ce^{C'(J+1)-cl}$.

The omitted atypical paths also have small \emph{relative} weight.
Indeed their fraction of the full mass is their probability under
the full tilt, and their fraction of the thresholded mass is their
probability under the thresholded tilt. Both are $Ce^{-cJ}$.
The factorization is used only on typical paths, where the already
charged counts agree; an overshoot below $h$ does not require such
a factorization. Taking $J$ a sufficiently small positive fraction
of $l$ proves \eqref{eq:radial-truncation} and completes the proof
of Proposition~\ref{prop:radial}.

\subsection{Several marked points}

Fix $p\ge2$ labeled tile vertices $\mathbf u=(u_1,\ldots,u_p)$.
For $\varnothing\ne I\subsetneq\{1,\ldots,p\}$ let $N_I$ be the number of plane loops
enclosing exactly the points with labels in $I$. Assign relative
weights $\mathbf v=(v_I)$, with $v_\varnothing=v_{\{1,\ldots,p\}}=1$,
and define
\[
 C(\mathbf v;\mathbf u)
       =\E\prod_{\varnothing\ne I\subsetneq\{1,\ldots,p\}}
                       v_I^{N_I}.
\]
This is well defined even for complex weights. Every separating loop
intersects a fixed finite collection of tile paths joining the points;
only finitely many switch strands use those paths. Let $b_0\in(1,2)$
be a fixed continuity cutoff for the limiting loop law, and write
\[
 F_n(v)=\E v^{\#\{\ell:u\text{ inside }\ell, R_u(\ell)\le b_0n\}}.
\]
Translation and duality make this function independent of the choice
of tile vertex $u$. Rounding and color conventions are transported
by these symmetries.

\begin{proposition}[Normalized multipoint correlations]
\label{prop:normalization}
Suppose $p$ is fixed, $\max_j|u_j|\le Cn$, and
$a=\min_{i\ne j}|u_i-u_j|\ge r_0$. Around every parameter whose
singleton weights $v_{\{j\}}$ are strictly positive real numbers,
the normalized correlation
\begin{equation}\label{eq:normalized-correlations}
 \mathcal C_n(\mathbf v;\mathbf u)
   =\frac{C(\mathbf v;\mathbf u)}
          {\prod_{j=1}^p F_n(v_{\{j\}})}
\end{equation}
is holomorphic and bounded in a fixed complex neighborhood of the
singleton weights, uniformly when the other weights range over a
bounded set, with bound $C(n/a)^C$. For $a\ge cn$, if $u_j/n$ converge
to distinct points and the point parities are fixed,
\eqref{eq:normalized-correlations} has a limit locally uniformly in
these complex parameter neighborhoods.

When all nontrivial weights lie in a fixed compact subset of
$(0,\infty)$ and $a\ge cn$, the normalized correlations are bounded
above and below by positive constants. Their tilted laws are local
at each point in the following quantitative sense. Consider two
possibly different collections with bounded point counts, all weights
in the stated positive compact set, and sizes comparable to the
same $n$, satisfying the same fixed geometric bounds. Align one
chosen point in each by a lattice symmetry, transporting its color,
and suppose its singleton weight is the same in both. The two laws admit
a coupling with identical configuration inside a common surrounding
circuit containing $B(u,s)$, with error $C(s/n)^c$. For minimum
separation $a$ the analogous error is at most
$C(n/a)^C(s/a)^c$, whenever $r_0\le s\le a/C$.
\end{proposition}

\begin{proof}
Take disjoint patches $P_j$ centered at $u_j$, of radii a small fixed
multiple of $a$, and with disjoint larger collars. At a positive
singleton base point $v_j=v_{\{j\}}$, use the positive auxiliary
weight
\[
 U=\prod_j v_j^{M_j},
 \qquad M_j=\#\{\ell:u_j\text{ inside }\ell,\ \ell\subset P_j\},
 \qquad Z_U=\E U.
\]
Write $\Pp_U=U\Pp/Z_U$. The collar comparison and
Proposition~\ref{prop:radial} give
\begin{equation}\label{eq:radial-proxy-mass}
 Z_U\asymp\prod_jF_a(v_j).
\end{equation}
Patches can be round, with a bounded whole-tile rounding; the mass
comparison is unaffected by a fixed change in their radii.

Every weighted loop not assigned to these local nests can be charged
to a passage at patch scale or larger. To make the count explicit,
join the points by a fixed spanning tree of paths of total length
$O(n)$. A separating loop hits a tree path. If this contact occurs
inside a patch but the loop is not contained there, follow the loop
to its first departure through a patch-scale buffer. Otherwise choose
its nearest approach to a marked point and use a buffer at the
corresponding geometric scale. Its diameter is comparable to or
larger than that scale. A loop enclosing several points has such
a witness as well. A bounded number of buffers at each scale
suffices because $p$ is fixed; choose them where the distance to the
nearest point is comparable to the scale. Their overlap has a
bounded coloring. Thus the total witness count $X$ satisfies
\begin{equation}\label{eq:radial-macro-moments}
 \E_Ue^{tX}\le C_t(n/a)^{C_t},\qquad t<\infty.
\end{equation}
For the lowest buffers, move the local cutoff back a bounded number
of levels and use the local tilt bounds before restoring it by
H\"older. Higher buffers are separated from the local tilts.
Conditioning through the disjoint patch collars treats all $p$
patches at once. This proves \eqref{eq:radial-macro-moments} directly
from the band estimates; it applies to arbitrarily large excursions,
since the finite tree already provides their witnesses.

At the positive singleton base, the absolute value of the remaining
weight is bounded by $K^X$, for a fixed $K$, even if some nonsingleton
weights vanish. Equations~\eqref{eq:radial-proxy-mass} and
\eqref{eq:radial-macro-moments}, followed by scale comparison from
$a$ to $n$, prove the asserted real-base upper bound. If all weights
are bounded above and below positively, the remaining weight is
also bounded below by $K^{-X}$. Jensen's inequality and
\eqref{eq:radial-macro-moments} then give the lower comparison.

For the complex bound, in every patch search inward for the first
stopping circuit with the same clear collars as before. Under
$\Pp_U$, its waiting time $L_j$ and the number $D_j$ of local charged
loops completed before it have a joint exponential moment, uniformly
in scale. The proof of \eqref{eq:radial-jump-tail} applies with a
patch cutoff in place of the initial circuit: its ordinary crossing
and collar estimates are conditional through the same free space.
The argument also applies to lower-threshold local tilts. Conditional
on these stopping circuits $\gamma_j$, both the true correlation
and the auxiliary weight factor into their exterior factors and
the continuation masses $F(\gamma_j;v_j)$.

Write the perturbed singleton weight as $v_je^{w_j}$. Before the
stops it changes the local weight by at most
$\exp(\sum_j|w_j|D_j)$ in absolute value. After dividing by the
reference factors in \eqref{eq:normalized-correlations},
\eqref{eq:radial-complex-ratio} bounds the continuation corrections
by
\[
 \exp\left(C\sum_j|w_j|
       [1+L_j+\log(2+n/a)]\right).
\]
The remaining nonsingleton factors are bounded by $K^X$. Apply
H\"older, use every fixed moment in
\eqref{eq:radial-macro-moments}, and then choose $\max_j|w_j|$ small
enough for the joint exponential moments of $L_j,D_j$.
This proves the uniform complex bound, including exhaustion at the
terminal scale. Nonvanishing of the denominator is supplied by
Proposition~\ref{prop:radial}.

We now prove convergence, retaining $a\ge cn$. First let the
singleton weights be real and positive; other weights may be
arbitrary fixed complex numbers in their bounded set. Remove
singleton weighting from all loops of maximum radius below
$\epsilon n$ about their point, and truncate each reference $F_n$
by exactly this prescription. Denote the resulting normalized
quantity by $\mathcal C_{n,\epsilon}$. The preceding bounds imply
\begin{equation}\label{eq:radial-normalized-truncation}
 \lim_{\epsilon\downarrow0}\limsup_{n\to\infty}
 |\mathcal C_n-\mathcal C_{n,\epsilon}|=0.
\end{equation}
Here are the details ensuring cancellation of the small-loop masses.
The first-screen waiting times in all patches can be restricted to
any fixed large number $L_0$ of bands with arbitrarily small additive
error, before or after thresholding. Use their exponential tails
and H\"older with the bounded moments of the absolute exterior
factor; the uniform threshold version gives the same estimate for
the truncated expression. On the retained stops the screens have
size at least a fixed multiple of $nK^{-L_0}$. For sufficiently
small $\epsilon$, every completed exterior loop with nontrivial
weight is unchanged by thresholding: a loop enclosing the point
outside its screen encloses that screen, and every other exterior
factor is already macroscopic. The only changes are therefore in
the continuation masses. Their ratios to the reference $F_n$ agree
to $1+o(1)$ by \eqref{eq:radial-truncation}, uniformly at these
boundedly many levels. The integral of the absolute exterior factor
is bounded. First let $n\to\infty$, then $\epsilon\downarrow0$,
and finally $L_0\to\infty$. This proves
\eqref{eq:radial-normalized-truncation}.

For fixed $\epsilon>0$, the truncated expression has a limit by
the plane nested-loop convergence in Input~\ref{inp:loops}. We
specify the continuity and integrability requirements. Loops deciding
the truncated test have spherical diameter bounded below after
rescaling, and the cited topology retains their individual curves.
A fixed continuum point lies on no macroscopic loop almost surely:
RSW screening bounds the probability that a loop arriving from a
fixed distance reaches a shrinking disk about it. Enclosure is
therefore a continuous winding test at each prescribed point.
Choose maximum-radius thresholds outside the at most countable
set of atoms of the relevant radius variables. Loops separating
the marked points but escaping an arbitrarily large fixed ball
can first be discarded. They must meet the finite joining tree,
so successive annular screens make the discarded probability tend
to zero uniformly in the mesh. Passage moments give uniform
integrability, including for complex exterior factors. Thus
convergence of the truncated expectations follows from the stated
plane theorem, without a bounded-domain loop limit. The real
truncated reference masses stay bounded away from zero: their
counts have uniformly bounded moments at fixed $\epsilon$, and
Jensen's inequality applies when the weight is below one.

The cutoff $b_0$ can be chosen once for the finitely many limiting
reference tests. Plane homogeneity gives the same choice at
translated points. Equation~\eqref{eq:radial-normalized-truncation}
then gives convergence of the untruncated normalized expression
for real positive singletons. The complex bounds give a normal
family. Two subsequential holomorphic limits agree on the real
singleton box, with arbitrary fixed remaining weights; repeated
one-variable uniqueness makes them identical throughout each
complex neighborhood. Hence the convergence is locally uniform.

Finally assume all weights are positive and bounded above and
below. The true tilted mass compares both ways to $Z_U$ by the
$K^{\pm X}$ bounds. Under that true law, H\"older and
\eqref{eq:radial-macro-moments} preserve an exponential tail for
the first patch screen, with the additional factor $(n/a)^C$
when the separation varies. Conditional on that screen, the inward
law is exactly the singleton tilt at its point. Apply the circuit
coupling from Proposition~\ref{prop:radial}, allowing for its random
initial level by the first-screen tail. This yields the final
coupling assertion. It also couples local observables obtained by
averaging orientation marks on loops contained in the common screen:
a mark strictly inside the screen cannot lie on an exterior loop.
\end{proof}

\begin{corollary}[Macroscopic tests with local insertions]
\label{cor:macroscopic-test}
Let $u_j/n$ converge to distinct points, with their point parities
fixed, and let $P_j$ be disjoint round patches centered at $u_j$,
of radii fixed positive multiples of $n$, with disjoint larger
collars. Fix $v_j>0$ and define
\[
 U_n=\prod_jv_j^{\#\{\ell:u_j\text{ inside }\ell,\ \ell\subset P_j\}}.
\]
Suppose $T_n$ is a uniformly bounded discretization,
$\sup_n\|T_n\|_\infty<\infty$, of a fixed bounded functional $T$ of the
rescaled plane loop collection, determined in a fixed bounded region
by loops of diameter at least a fixed $\delta>0$. Suppose that this
functional is almost surely continuous for the plane loop limit,
with its marked-point, patch-boundary and contact tests taken at
continuity radii. Here discretization means that
$T_n(\omega_n)\to T(\omega)$ whenever the relevant rescaled
loop collections $\omega_n$ converge to a continuity configuration
$\omega$. Then the following sequence converges as $n\to\infty$:
\begin{equation}\label{eq:radial-macroscopic-test}
 \frac{\E[U_nT_n]}{\prod_jF_n(v_j)}.
\end{equation}

The conclusion also applies to tests obtained by first excluding
arbitrarily large excursions, provided the discarded part tends
to zero in probability under the normalized $U_n$ law uniformly
in $n$. Boundedness of $T_n$ may be replaced by
$|T_n|\le K^{X_n}$, where $X_n$ is the sum of a fixed finite family of
macroscopic passage counts satisfying
\begin{equation}\label{eq:radial-test-integrability}
 \sup_n\frac{\E[U_ne^{tX_n}]}{\E U_n}<\infty
 \qquad\text{for every fixed }t<\infty,
\end{equation}
also after the inner-radius truncations of the local weights.
The bounded truncations of the test must satisfy the preceding
continuity requirement.
\end{corollary}

\begin{proof}
Choose $\eta>0$ smaller than the patch radii and, taking account of
the fixed band-radius factors, small enough that every screen sought
below $\eta n$ has diameter less than $\delta n$.
Under the normalized $U_n$ law, regeneration between radii
$\eta n$ and $\epsilon n$ supplies a surrounding circuit at each
marked point, except with probability $C(\epsilon/\eta)^c$.
The same statement holds for thresholded local weights. A loop
used by the test has diameter at least $\delta n$, so cannot be
contained inside these circuits; it cannot cross them either.
Thus the events where a tested macroscopic loop enters the
$\epsilon n$ neighborhoods of the marked points can be discarded
with an error tending to zero. Under
\eqref{eq:radial-test-integrability}, H\"older's inequality makes
the same discarded contribution small for the unbounded test.

Search for the first screens below $\eta n$. Conditional on these
screens the test, after this exclusion, is determined by their
exterior. Restrict their waiting times to a fixed number of bands,
use their exponential tails for the remainder, and truncate the
local weights and all reference factors at the same radius
$\epsilon n$. Exactly the continuation-quotient argument proving
\eqref{eq:radial-normalized-truncation} shows that this changes
\eqref{eq:radial-macroscopic-test} by $o(1)$, uniformly for large
$n$ as $\epsilon\downarrow0$. For fixed $\epsilon$, all remaining
weights and tests concern macroscopic loops. Input~\ref{inp:loops}
and the stated continuity assumptions give convergence. Large
excursions are removed before this step and restored afterwards
using their assumed small probability and uniform integrability.

For an unbounded test, truncate at $|T_n|\le M$. The mass comparison
\eqref{eq:radial-proxy-mass} bounds
$\E U_n/\prod_jF_n(v_j)$ above and below. Condition
\eqref{eq:radial-test-integrability} therefore makes the discarded
normalized expectations tend to zero as $M\to\infty$, uniformly
in $n$ and in the inner cutoff. This proves the extension.
For the passage counts used below, condition
\eqref{eq:radial-test-integrability} follows from
\eqref{eq:radial-macro-moments}: their buffers are outside the local
patch cores, with a bounded cutoff rollback at the innermost buffers.
\end{proof}
